\section*{الفصل \ref{ch:b2:unicellular-diversity} --- تنوع الكائنات وحيدة الخلية}

\begin{solution}{exo:b2:unicellular-diversity:1}
\omterm{def:b2:unicellular-diversity:unicellular}{وحيد الخلية}: خلية واحدة تؤدي كل الوظائف وتعيش وحدها؛ ومستعمري:
خلايا متماثلة تبقى ملتصقة لكن كل واحدة منها تستطيع العيش وحدها؛
\omterm{def:b2:unicellular-diversity:unicellular}{وعديد الخلايا}: أنواع خلوية عدة يعتمد بعضها على بعض، ولا يتكاثر منها
إلا السلالة الجنسية. و\emph{E.~coli} والخميرة وحيدا الخلية (البرعم
ينفصل ويعيش وحده)؛ و\emph{Gonium} مستعمري (16 خلية متماثلة، كل
منها قادرة على تأسيس مستعمرة)؛ و\emph{Volvox} \omterm{def:b2:unicellular-diversity:unicellular}{عديد الخلايا} (خلايا
جسدية فانية، وقليل من الخلايا الإنتاشية).
\end{solution}

\begin{solution}{exo:b2:unicellular-diversity:2}
$S/V = 3/r$: أي \qty{6}{\micro m^{-1}} عند المكوّر،
و\qty{0.06}{\micro m^{-1}} عند \omterm{def:b2:unicellular-diversity:protist}{الطلائعي}، بفرق مئة ضعف. والمكوّر
يتنفس أسرع لكل وحدة حجم: فكل جزء من سيتوبلازمه يقع على بعد جزء من
الميكرومتر من السطح الذي يدخل عبره الأكسجين والغذاء، بينما يُغذَّى
داخل \omterm{def:b2:unicellular-diversity:protist}{الطلائعي} عبر سطح أقل مئة مرة لكل وحدة حجم.
\end{solution}

\begin{solution}{exo:b2:unicellular-diversity:3}
لا نواة (بل \omterm{def:b2:unicellular-diversity:prokaryote}{جسم نووي}) في مقابل غلاف نووي؛ ولا متقدرات ولا أغشية
داخلية، وسلسلة التنفس في الغشاء الهيولي، في مقابل عضيات؛ وسوط هو
حلزون صلب من الفلاجيلين تديره البروتونات، في مقابل هدب $9+2$
ينثني بالداينين وبجزيئات ATP؛ وكذلك القياس (\qty{1}{\micro m} في
مقابل عشرات) وجدار من الببتيدوغليكان. ولا تستطيع البكتيريا البلعمة:
جدارها يمنع ابتلاع جسيم، وهو ما تفعله الأميبا بأقدامها الكاذبة.
\end{solution}

\begin{solution}{exo:b2:unicellular-diversity:4}
الأهداب: الحركة وكنس الغذاء نحو الفم؛ والميزاب الفموي: يقمع
الجسيمات إلى البلعوم؛ والفجوة الغذائية: تهضم البكتيريا المبتلعة
بإنزيمات جسيمية حالّة؛ \omterm{prop:b2:unicellular-diversity:osmoregulation}{والفجوة المتقلصة}: تطرد الماء الداخل
بالتناضح؛ والنواة الكبيرة: النواة العاملة، تُستنسخ منها مئات نسخ
المورّثات لحياة كل يوم؛ والنواة الصغيرة: النواة الإنتاشية الثنائية
الصيغة، تُستعمل في التقسم المنصّف والاقتران.
\end{solution}

\begin{solution}{exo:b2:unicellular-diversity:5}
$t = x^2/2D$: أي $(10^{-6})^2/10^{-9} = \qty{1}{ms}$ عند البكتيريا؛
و$(5\times 10^{-4})^2/10^{-9} = \qty{250}{s}$، أي نحو أربع دقائق،
عند الأميبا. ولا تستطيع الأميبا الاعتماد على الانتشار في توزيع
المستقلَبات: فهي تحتاج إلى جريان سيتوبلازمي ونقل تدفعه محركات على
امتداد الهيكل الخلوي، وهو ما تستغني عنه البكتيريا.
\end{solution}

\begin{solution}{exo:b2:unicellular-diversity:6}
$Re = \rho vL/\eta$: عند \emph{Paramecium} $10^3\times 10^{-3}\times
2\times 10^{-4}/10^{-3} = 0.2$؛ وعند بكتيريا $10^3\times 2\times
10^{-5}\times 2\times 10^{-6}/10^{-3} = 4\times 10^{-5}$؛ وعند سابح
$10^3\times 1\times 2/10^{-3} = 2\times 10^6$. وعند $Re \gg 1$ تحمل
العطالة السابح بعد الضربة؛ وعند $Re \ll 1$ توقف قوة الجرّ اللزجة
البكتيريا خلال نانومتر، فلا تتحرك إلا ما دامت تضرب.
\end{solution}

\begin{solution}{exo:b2:unicellular-diversity:7}
$v_s = 2r^2\Delta\rho g/9\eta = 2\times 10^{-10}\times 100\times
9.81/(9\times 10^{-3}) = \qty{2.2e-5}{m/s}$؛ ويستغرق \qty{50}{m}
$50/2.2\times 10^{-5} = \qty{2.3e6}{s}$، أي 27 يومًا. وتنصيف نصف
القطر يقسم $v_s$ على 4: 108 أيام. وخفض $\Delta\rho$ إلى
\qty{20}{kg/m^3} يقسمها على 5: 135 يومًا.
\end{solution}

\begin{solution}{exo:b2:unicellular-diversity:8}
$V = \tfrac{4}{3}\pi\times 100\times 25\times 25 =
\qty{2.6e5}{\micro m^3}$. وفي \qty{1800}{s}: $2.6\times 10^5/1800 =
\qty{145}{\micro m^3/s}$؛ ولمّا كان $\qty{1}{\micro m^3} =
\qty{1e-15}{L}$، فهذا \qty{1.5e-13}{L/s}.
\end{solution}

\begin{solution}{exo:b2:unicellular-diversity:9}
التضاعف على امتداد \qty{1}{mm} ارتفاع قدره نحو \qty{0.07}{\%} لكل
ميكرومتر (لأن $2^{1/1000} - 1 = 0.0007$)، أي \qty{0.07}{\%} عبر
الخلية. والشوط الذي يدوم \qty{1}{s} يقطع \qty{20}{\micro m}: أي
\qty{1.4}{\%}. ولتمييز فرق قدره \qty{0.07}{\%} كان على الخلية أن
تعدّ نحو $(1/0.0007)^2 = 2\times 10^6$ جزيئة عند كل طرف، وهو أكثر
بكثير مما ترتبط به مستقبلاتها التي لا تتجاوز بضعة آلاف في الثانية؛
أما \qty{1.4}{\%} فتحتاج نحو \num{5000}، وهذا ممكن. فالمقارنة
الزمنية تحوّل قياسًا فراغيًا ميئوسًا منه إلى قياس ممكن، إذ تدع
الشوط يقوم بالمكاملة.
\end{solution}

\begin{solution}{exo:b2:unicellular-diversity:10}
السطح $4\pi r^2$ مع $r = \qty{250}{\micro m}$: أي \qty{7.9e5}{\micro
m^2}؛ والحجم السيتوبلازمي $\approx$ السطح $\times$ \qty{2}{\micro m}
$= \qty{1.6e6}{\micro m^3}$؛ فالنسبة \qty{0.5}{\micro m^{-1}}. أما
\emph{E.~coli}، وهي أسطوانة: $S/V = 2\pi r L/\pi r^2 L = 2/r =
\qty{4}{\micro m^{-1}}$. فالعملاقة أسوأ من \emph{E.~coli} بثماني
مرات فقط لا بخمسمئة مرة، لأن كل نقطة من سيتوبلازمها تقع على بعد
\qty{2}{\micro m} من السطح. وتخزن الفجوة النترات، وهي متقبل
إلكتروناتها، لتستطيع تنفس الكبريتيد في الرسابة اللاأكسجينية بين
التقليبات النادرة التي تجلب النترات
(\cref{ch:b2:microbial-metabolism}).
\end{solution}

\begin{solution}{exo:b2:unicellular-diversity:11}
لا تستطيع الخلية أن تسبح وتنقسم في آن. فمستعمرة من 16 خلية تكفّ عن
السباحة ساعات الانقسام تهبط، \omterm{thm:b2:unicellular-diversity:stokes}{بقانون ستوكس}، مسافة مهملة (نصف قطرها
بضع عشرات من الميكرومترات، وسرعة هبوطها ميكرومترات في الثانية)؛
فتستطيع المستعمرة كلها أن تنقسم دون خسارة تُذكر. أما كرة من
\num{10000} خلية نصف قطرها \qty{500}{\micro m} فكانت لتهبط أسرع
بمئة مرة --- مليمترات في الثانية، وأمتارًا في الساعة --- خارج
الضوء: فيجدي إبقاء معظم الخلايا سابحة بينما ينقسم قليل منها.
والخلايا الإنتاشية كبيرة لأن عليها أن تحمل الموارد اللازمة لبناء
كرة بنت كاملة من آلاف الخلايا بانقسامات متتابعة دون تغذٍّ، فتمونها
الخلايا الجسدية.
\end{solution}

\begin{solution}{exo:b2:unicellular-diversity:12}
الفرع الحيوي سلف وكل ذريته؛ أما الدرجة فمستوى من التنظيم بلغته
سلالات عدة. \omterm{def:b2:unicellular-diversity:protist}{والطلائعيات} --- الطحالب الخضراء (مع النباتات اليابسة)،
والدياتومات (مع الطحالب البنية)، والهدبيات (مع السوطيات الدوّارة
و\emph{Plasmodium})، والأميبات (قرب الحيوانات والفطريات) --- لا
تشترك إلا في الخلية حقيقية النواة تعيش وحدها، وسلفها المشترك الأخير
هو سلف حقيقيات النوى كلها، ونحن منها: فهي درجة. و«بدائيات النوى»
تجمع البكتيريا والعتائق بما تفتقر إليه؛ والعتائق أقرب إلى حقيقيات
النوى منها إلى البكتيريا، فالمصطلح درجة هو أيضًا. أما الزراقم
فتنحدر من سلف واحد ابتكر التركيب الضوئي المنتج للأكسجين، وتضم
ذريته كلها (باستثناء سلالة الصانع الأخضر): فهي فرع حيوي.
\end{solution}

\begin{solution}{pb:b2:unicellular-diversity:1}
\textbf{1.} $1\times 1\times 0.1 = \qty{0.1}{mm^3} = \qty{0.1}{\micro L}$.
\textbf{2.} $24/\qty{0.1}{\micro L} = 240$ لكل \unit{\micro L} $=
\qty{2.4e5}{}$ خلية في المليلتر.
\textbf{3.} $V = \tfrac{4}{3}\pi\times 125 = \qty{524}{\micro m^3}$;
$S/V = 3/5 = \qty{0.6}{\micro m^{-1}}$.
\textbf{4.} $2.4\times 10^5\times 524 = \qty{1.26e8}{\micro m^3}$ في
المليلتر، و$\qty{1}{mL} = \qty{1e12}{\micro m^3}$: أي كسر قدره
$1.3\times 10^{-4}$، \qty{0.013}{\%}.
\textbf{5.} $\qty{1000}{m^3} = \qty{1e9}{mL}$: $2.4\times 10^{14}$
خلية.
\textbf{6.} الزراق خيط من خلايا أصغر بكثير (بضعة ميكرومترات)،
خضراء مزرقّة بلا صانع أخضر متميز ولا نواة، وهو لا يسبح بأسواط؛ أما
الطحلب فخلية بيضوية مفردة لها سوطان وبقعة عينية وصانع أخضر واحد على
شكل كأس.
\textbf{7.} $384/24 = 16 = 2^4$: أي أربع مضاعفات في \qty{48}{h}،
$T = \qty{12}{h}$.
\textbf{8.} $k = \ln 2/T = 0.693/12 = \qty{0.058}{h^{-1}}$.
\textbf{9.} $t = \ln(10^8/2.4\times 10^5)/k = \ln(417)/0.058 =
6.03/0.058 = \qty{104}{h}$، أي نحو 4.3 يوم.
\textbf{10.} تنفد المغذيات (الفوسفات والنترات)؛ وتظلّل الخلايا
بعضها بعضًا فيهبط نصيب الخلية من الضوء؛ وتتكاثر عليها الراعيات
(الهدبيات، والدولابيات، و\emph{Daphnia})؛ وعند \qty{1e8}{} خلية في
المليلتر تملأ الخلايا نحو خمسة في المئة من الحجم وتستنفد
$\mathrm{CO_2}$ المنحلّ.
\textbf{11.} النمو نصف الوقت فقط: فينتصف المعدل المتوسط،
\qty{208}{h}، أي 8.7 يوم.
\textbf{12.} يقيس العدّ الجمهرة لا الخلايا فرادى: فإذا أطلقت خلية 8
بنات كل \qty{36}{h}، ضُربت الجمهرة في 8 $= 2^3$ خلال ثلاثة أزمنة
جيل مقدار كل منها \qty{12}{h}. \omterm{def:b2:unicellular-diversity:division}{فزمن الجيل} هو متوسط زمن تضاعف العدد،
مهما كانت طريقة تجميع الانقسامات.
\textbf{13.} $Re = 10^3\times 10^{-4}\times 10^{-5}/10^{-3} = 10^{-3}$.
\textbf{14.} $m = 1050\times 5.24\times 10^{-16} = \qty{5.5e-13}{kg}$؛
و$6\pi\eta r = 6\pi\times 10^{-3}\times 5\times 10^{-6} =
\qty{9.4e-8}{kg/s}$؛ ومنه $d = 5.5\times 10^{-13}\times 10^{-4}/9.4\times
10^{-8} = \qty{5.8e-10}{m}$: أي نصف نانومتر.
\textbf{15.} $v_s = 2\times 25\times 10^{-12}\times 50\times
9.81/(9\times 10^{-3}) = \qty{2.7e-6}{m/s}$، أي \qty{2.7}{\micro m/s}.
\textbf{16.} الهبوط \qty{1}{m}: $1/2.7\times 10^{-6} =
\qty{3.7e5}{s}$، أي 4.3 يوم؛ والسباحة صعودًا بسرعة \qty{100}{\micro m/s}:
\qty{1e4}{s}، أي 2.8 ساعة. فالسباحة تكسب بعامل 37.
\textbf{17.} $c_\infty = \qty{1e-2}{mol/m^3}$: $J = 4\pi\times
1.9\times 10^{-9}\times 5\times 10^{-6}\times 10^{-2} =
\qty{1.2e-15}{mol/s} = 7.2\times 10^8$ جزيئة في الثانية.
\textbf{18.} \qty{100}{pg} $= \qty{1e-10}{g} = \qty{8.3e-12}{mol}$ من
الكربون، تتضاعف في \qty{12}{h} $= \qty{43200}{s}$:
\qty{1.9e-16}{mol/s} $= 1.2\times 10^8$ ذرة في الثانية. ونسبة الوارد
إلى الطلب: $7.2/1.2 = 6$.
\textbf{19.} ينمو الوارد عشرة أضعاف ($\propto r$)، وينمو الطلب ألف
ضعف ($\propto r^3$): فتصير النسبة $6/100 = 0.06$؛ وخلية كهذه لا
تستطيع النمو إلا بستة في المئة من المعدل، فطحلب بذلك القياس يتغذى
بالانتشار مستحيل دون تقليب أو نقل داخلي أو استقلاب أبطأ بكثير.
\textbf{20.} $\Delta\Pi = RT\Delta c = 8.314\times 293\times 99 =
\qty{2.4e5}{Pa}$، أي \qty{2.4}{bar}.
\textbf{21.} المساحة $4\pi(5\times 10^{-6})^2 = \qty{3.14e-10}{m^2}$؛
والوارد $L_p A\Delta\Pi = 10^{-14}\times 3.14\times 10^{-10}\times
2.4\times 10^5 = \qty{7.6e-19}{m^3/s} = \qty{0.76}{\micro m^3/s}$.
\textbf{22.} $524/0.76 = \qty{690}{s}$، أي نحو إحدى عشرة دقيقة.
\textbf{23.} $P = \Delta\Pi\times$ التدفق $= 2.4\times 10^5\times
7.6\times 10^{-19} = \qty{1.8e-13}{W}$.
\textbf{24.} جزيء ATP واحد: $5\times 10^4/6.02\times 10^{23} =
\qty{8.3e-20}{J}$: ومنه $1.8\times 10^{-13}/8.3\times 10^{-20} = 2.2\times
10^6$ من ATP في الثانية. أما تثبيت الكربون فيكلف $3\times 1.2\times 10^8 =
3.6\times 10^8$ من ATP في الثانية. فالفجوة تكلف \qty{0.6}{\%} من
ذلك --- وهو تأمين رخيص ضد الانفجار.
\textbf{25.} وارد الماء \qty{0.76}{\micro m^3/s} (أي حجم الخلية
في إحدى عشرة دقيقة)؛ ولولا \omterm{prop:b2:unicellular-diversity:osmoregulation}{فجوة متقلصة} لضاعفت الخلية حجمها في نحو
\qty{700}{s}؛ وبقاؤها جافة يكلف على الأقل
$2\times 10^6$ من ATP في الثانية، أي أقل من واحد في المئة مما ينفقه
تركيبها الضوئي على الكربون.
\end{solution}
